\begin{abstract}

\vspace{0.6em}
\noindent\textbf{摘要}——监控摄像头视频中的异常事件检测是公共安全的核心技术，但走向真实部署时面临三重困难：通用视觉-语言大模型与监控域之间存在显著域偏移，异常分数整体偏低；异常解释由大模型自由生成，无法判断其中哪些内容真正有视频证据支撑；免训练流程隐含假设输入足够干净，低光、雨雾、视频压缩与相机抖动会使其评分性能明显下降并产生幻觉语义。本文提出 \evivad{}，一个在完全冻结的视觉-语言大模型之上构建的统一框架，包含三个互补模块：（1）\textbf{DAA}（Domain Low-rank Adapter）在视觉编码器末端四层的 Q/V 投影注入秩为 $r{=}4$ 的低秩增量，以不足 $5$\,M 的可训练参数完成监控域适配；（2）\textbf{EAD}（Evidence-Anchored Decoding）强制解释显式引用时间区间与视觉线索，并以 InfoNCE 形式的证据对齐损失进行训练，使解释既可被证据支持、也可被反例推翻；（3）\textbf{DAG}（Degradation-Aware Gating）以轻量质量探针调制「视觉-文本对齐」与「检索/先验主导」两条打分路径，并在低置信时显式弃权而非给出高置信的错误告警。我们在 UCF-Crime 主数据集与 XD-Violence、UBnormal、MSAD 三个跨域集上，以及在 $4$ 类退化 $\times$ $3$ 个强度共 $12$ 种退化配置下建立了统一评测协议。结果显示：相对免训练基线 B0，\evivad{} 将帧级 AUC 由 $76.8$ 提升至 $82.1$，证据归因召回（EAR）由 $41.5$ 提升至 $68.4$，幻觉率（HR）由 $26.4$ 降至 $9.8$，退化网格平均 AUC 由 $62.5$ 提升至 $71.3$（相对性能保持率由 $0.81$ 升至 $0.92$），而可训练参数仅增加 $5.5$\,M。本文的主要意义在于指出：大模型时代的视频异常检测评价不应只看检测分数，\emph{解释的可验证性}与\emph{输入退化下的可降级性}应当成为同等重要的核心评估指标。

\vspace{0.6em}
\noindent\textbf{Abstract}--- Anomaly detection in surveillance camera videos is a cornerstone of public safety, yet real deployments face three obstacles: general vision-language models suffer a severe domain shift from the surveillance domain; free-form textual explanations cannot be verified against the video evidence; and training-free pipelines implicitly assume clean inputs, so low light, rain/fog, compression and camera shake collapse their scores and induce hallucinated semantics. We present \evivad{}, a unified framework built on a \emph{fully frozen} vision-language backbone with three complementary modules: (1) \textbf{DAA} injects rank-$4$ low-rank increments into the Q/V projections of the last four encoder blocks, adapting to the surveillance domain with fewer than $5$\,M trainable parameters; (2) \textbf{EAD} forces explanations to explicitly cite temporal intervals and visual cues, trained with an InfoNCE evidence-alignment loss so that an explanation can be both supported by evidence and falsified by counterexamples; (3) \textbf{DAG} uses a lightweight quality probe to blend a vision-text alignment path with a retrieval/prior-driven path, and abstains under low confidence instead of emitting confident false alarms. Under a unified protocol on UCF-Crime and three cross-domain sets (XD-Violence, UBnormal, MSAD) plus a $4\times3$ degradation grid, The aggregate results indicate that \evivad{} improves frame-level AUC from $76.8$ to $82.1$, evidence-attribution recall from $41.5$ to $68.4$, reduces the hallucination rate from $26.4$ to $9.8$, and raises the average AUC on the degradation grid from $62.5$ to $71.3$ (relative performance retention $0.81 \rightarrow 0.92$) with only $5.5$\,M additional trainable parameters. Our main message is that in the era of large models, video anomaly detection should be evaluated not by detection scores alone: the \emph{verifiability} of explanations and the \emph{graceful degradability} under input corruption deserve equal standing.
\end{abstract}
